\section{Saturation removes an avoidable rank loss}\label{sec:saturation}

\status{This strengthens the local accounting for the actual retained reference lattice. The joint exceptional-direction estimates remain inputs. The reference-lattice comparisons and descent used below are explicit in the pinned source; the elementary deductions are proved here.}

Fix a word $u$ and a served finite place $w\mid p$, so $p\ne T_u$.
Write $K$ for the coefficient field, $A=\mathcal O_{K,w}$, and
$\mathcal O\subset F$ for the source's comparison DVR and its fraction field.
Let $L_u$ be the full integral column lattice, of rank $d_u$, and let
$U_u$ be its rational ramified subspace. Its retained reference lattice is
$R_u=L_u\cap U_u$, of rank $r_u>0$.
The source's coordinate inclusion identifies this with \texttt{DV.MramL}.
At $w$, the comparison map $\phi$ identifies
\begin{equation}\label{eq:sr-reference}
 M=\mathcal O\phi(L_u),\qquad
 U_F=F\phi(U_u),\qquad
 M_{\rm ram}=M\cap U_F=\mathcal O\phi(R_u).
\end{equation}
Here $M$ is \texttt{G.M}, not \texttt{G.Mna}. Since $U_F$ is a vector
subspace, $ax\in U_F$ and $a\ne0$ imply $x\in U_F$. Thus
$M_{\rm ram}$ is saturated in $M$, and
\[
 M_{\rm ram}\cap p^tM=p^tM_{\rm ram}\qquad(t\ge0).
\]
The same argument proves saturation of $R_u$ in $L_u$ before comparison.
No tensor-model or metric identification is needed for this fact.
The selected finite-index lattice need not itself be saturated. \textsf{[SR1]}

\paragraph{The actual comparison and descent.}
The source specifies an embedding $\sigma:K\hookrightarrow F$ with
$\sigma(x)\in\mathcal O$ if and only if $x\in A$, and requires images
of a $K$-basis under $\phi$ to remain $F$-linearly independent.
Consequently $A\to\mathcal O$ is an injective local map of DVRs.
It is faithfully flat, and scalar extension commutes with intersections
of lattices and rational subspaces. This proves \eqref{eq:sr-reference}.
The actual local selected lattice is
\[
 P=J\cap p^\tau M\cap U_F,
 \quad
 J=\sum_{B\subseteq\mathcal V}
 p^{\sum_{v\notin B}S_v}\Bigl(\prod_{v\in B}y_v^m\Bigr)M.
\]
The recorded rationality of the $y_v$ makes $J$ an extension of an
$A$-lattice; allowing a unit factor for each $y_v$ gives the same conclusion.
Flatness then shows that $P$ is an extension of its rational preimage.
The source's served-place identity identifies that preimage with
$(N_u)_w$, where $N_u=\texttt{NuLat}\,u$.
This is the content of \texttt{M\_inf\_Vram}, \texttt{NP\_eq\_span}, and
\texttt{length\_quotient\_of\_nonempty}, with rationality supplied through
\texttt{NPDesc\_of}. It is not an assumption that two commensurable
reference lattices have the same index.

Put $e_w=v_w(p)$ and $e_{\mathcal O}=v_{\mathcal O}(p)$.
If $a=v_{\mathcal O}(\pi_w)$, then extension multiplies each Smith
exponent by $a$, and $e_{\mathcal O}=ae_w$. Hence
\begin{equation}\label{eq:sr-length-transport}
 \frac{\operatorname{length}_{A}((R_u)_w/(N_u)_w)}{e_w}
 =\frac{\operatorname{length}_{\mathcal O}(M_{\rm ram}/P)}{e_{\mathcal O}}.
\end{equation}
This also preserves the ranks $d_u,r_u$. \textsf{[SR2]}

\paragraph{Why the smaller full lattice causes only a depth shift.}
The comparison lattice $M_{\rm na}=\texttt{G.Mna}$ satisfies
$p^eM_{\rm na}\subseteq M\subseteq M_{\rm na}$.
Both are preserved by the $y_v$. If $J_{\rm na}$ is defined as $J$ using
$M_{\rm na}$, then $p^eJ_{\rm na}\subseteq J$.
Thus, from
$p^EM_{\rm na}\subseteq J_{\rm na}+\sum_{i=1}^g\mathcal O z_i$,
multiplication by $p^e$ gives
\[
 p^{E+e}M\subseteq J+\sum_{i=1}^g\mathcal O(p^ez_i),
 \qquad p^ez_i\in M.
\]
The exceptional count is unchanged, and both full lattices have the same
rank. This is precisely the source's \texttt{inverse\_shift} step.
For a nonempty served set the resulting depth is the integer
\[
 \tau=e+\sum_{v\in\mathcal V}m(t_v+D_0+1).
\]
The real-valued \texttt{tauP} is this integer, and \texttt{Sv} is already
a nonnegative integer. Their displayed ceilings introduce no additional
rounding loss. For an empty served set, $\tau=\mathrm{io}=0$ and
$P=M_{\rm ram}$. At unserved regular places the recorded
\texttt{NuUnserved} identity gives $(N_u)_w=(R_u)_w$, hence zero index.
\textsf{[SR3]}

\begin{proposition}\label{prop:sr-actual-index}
Assume the source's joint exceptional-direction bounds: at depth
$\tau+\beta S$, $J\subseteq M$ has at most
$d_u m^{\beta+1}e_{\beta+1}(\kappa)$ exceptional directions for every
$\beta\ge0$, with $S\ge S_v$, $\kappa_v\ge0$ and
$\Sigma=m\sum_v\kappa_v\le1/2$.
Set $\Theta=\prod_v(1+m\kappa_v)-1$ and $\mathrm{io}=2S\Sigma$.
Then the actual selected lattice satisfies
\begin{equation}\label{eq:sr-index}
 \operatorname{length}_{A}((R_u)_w/(N_u)_w)
 \le e_w\{r_u\tau+d_uS\Theta\}
 \le e_w\{r_u\tau+d_u\mathrm{io}\}.
\end{equation}
\end{proposition}

\begin{proof}
Put $Q=J\cap p^\tau M$. Saturation gives an exact sequence
\[
 0\longrightarrow p^\tau M_{\rm ram}/P
 \longrightarrow M_{\rm ram}/P
 \longrightarrow M_{\rm ram}/p^\tau M_{\rm ram}\longrightarrow0.
\]
The last quotient has length $e_{\mathcal O}r_u\tau$.
The first embeds in $p^\tau M/Q$, whose length is at most
$e_{\mathcal O}d_uS\Theta$ by the elementary-divisor tail sum in LC2.
Apply \eqref{eq:sr-length-transport}, then $\Theta\le2\Sigma$.
\end{proof}

Writing $\rho_u=d_u/r_u$, the slope conclusion of LC5 therefore improves to
\begin{equation}\label{eq:sr-slope}
 \mu(I_uN_u)\ge\mu(R_u)+\mathrm{Ga}+\mathrm{Fo}-\mathrm{Ee}_u
 -\rho_u\mathrm{IO}-\mathcal R_u-\mathcal B_u.
\end{equation}
The entire $(\rho_u-1)\mathrm{Ta}$ term is removed. The type prime
$T_u$ still has its separately estimated cost $\mathcal B_u$.
The actual source sizes give $T_u\ge X'$, so its served set is empty and
$\tau=\mathrm{io}=0$ there. Thus the exceptional baseline in LC5 is also
zero, and $\mathcal B_u$ retains its previous meaning. If any additional
exceptional places were introduced with nonzero baseline, their excess
would have to be estimated relative to the new baseline. \textsf{[SR4]}

\paragraph{The source's rank sizes absorb the remaining ratio.}
One can retain more of the same tail estimate:
\[
 \Theta\le\sum_{j\ge1}\frac{\Sigma^j}{j!}
 \le\frac{\Sigma}{1-\Sigma/2}\le\frac43\Sigma,
\]
because $j!\ge2^{j-1}$. If $r_u/d_u\ge2/3$, then
$d_uS\Theta\le r_u(2S\Sigma)=r_u\mathrm{io}$.
The source's actual rank formula and size hypotheses give this with room
to spare:
\[
 \frac{r_u}{d_u}=\left(1-\frac1{T_u-1}\right)^f,
 \quad16f\le b,\quad4\le b,\quad b^3\le T_u
 \ \Longrightarrow\ 
 1-\frac{r_u}{d_u}\le\frac{2f}{T_u}
 \le\frac1{8b^2}\le\frac1{128}.
\]
These numerical hypotheses already occur in \texttt{PlacesP.sizes\_at}
and its family assembly. Thus, under the same retained arithmetic and
comparison inputs, \eqref{eq:sr-slope} strengthens to
\begin{equation}\label{eq:sr-final}
 \boxed{\mu(I_uN_u)\ge\mu(R_u)+\mathrm{Ga}+\mathrm{Fo}-\mathrm{Ee}_u
 -\mathrm{IO}-\mathcal R_u-\mathcal B_u.}
\end{equation}
If the actual integral and isometric transport required in LC7 is also
supplied, its separate metric identity is
$\mu(R_u)=\mu(L_u)-f\log(T_u-1)/(T_u-2)$.
The mean correction, the finite-place rounding, and the type-prime cost
remain explicit. In the rational-prime notation $\mathrm{IO}=2\mathrm{IOp}$,
the regular-place allowance is now $2\mathrm{IOp}$ instead of the source's
$4\mathrm{IOp}$ after its coarser absorption. \textsf{[SR5]}
